\chapter{Fundamentals}\label{ch:fundamentals}

This chapter establishes the background needed for the remainder of the thesis: the
cache hierarchy and the notion of cache awareness, trie-based index structures with
the Adaptive Radix Tree as the running reference, the unified \texttt{std::map}
comparison interface, the YCSB workloads, and the C++23 metaprogramming facilities on
which the cache engine is built.

\section{Cache Hierarchy and Cache Awareness}\label{sec:cache-basics}

Modern processors interpose several levels of cache (L1, L2, L3) between the registers
and main memory. Data is transferred between these levels in fixed units called
\emph{cache lines}, typically 64\,bytes. An access that finds its datum in a cache
level is cheap; a miss propagates to the next, slower level, and a miss in the last-level
cache reaches main memory at a cost of hundreds of cycles. Address translation adds a
second hierarchy: the \emph{translation lookaside buffer} (TLB) caches virtual-to-physical
mappings, and a dTLB miss is itself expensive. On multi-socket systems, non-uniform memory
access (NUMA) makes the latency of a reference depend on which node owns the page.

A data structure is \emph{cache-aware} when its layout is tuned to concrete cache
parameters (line size, level capacities), and \emph{cache-oblivious} when it achieves good
locality across all levels without knowing them. Both philosophies recur throughout the
state of the art (Chapter~\ref{ch:sota}) and motivate the layout, prefetch, and hardware
axes of the proposed engine.

\section{Trie-Based Index Structures}\label{sec:trie-basics}

A \emph{trie} (digital/radix tree) indexes keys by their byte sequence rather than by
comparison: each level discriminates on one or more key bytes, so the search-path length
depends on key length, not on the number of stored keys. The \emph{fanout} (branching
factor) and the \emph{node representation} dominate both space and cache behaviour. The
Adaptive Radix Tree (ART)~\cite{leis2013art} adapts the node representation to the actual
fanout, using four node types (Node4, Node16, Node48, Node256), and applies path
compression (lazy expansion, prefix folding) to keep the tree shallow. ART serves as the
running reference design throughout this thesis.

\section{The Unified Comparison Interface (\texorpdfstring{\texttt{std::map}}{std::map})}\label{sec:stdmap-interface}

To compare structurally different search algorithms fairly, they are all mapped onto a
single, well-understood interface: the C++ ordered-map contract
\texttt{std::map<Key,\,Value>} (\texttt{find}/\texttt{insert}/\texttt{erase}/range
iteration). The axes introduced in Chapter~\ref{ch:concept} describe the \emph{internal}
behaviour behind this interface, while the interface itself stays fixed --- comparing two
conforming \texttt{std::map} implementations is therefore the central measurement act. A
variadic specialization keeps the interface ergonomic: one template parameter exposes a
\texttt{std::vector}-like API, two parameters the \texttt{std::map} API, and $N>2$ a
\texttt{std::map<K,\,std::tuple<\ldots>>}.

\section{YCSB Workloads}\label{sec:ycsb}

The Yahoo!\ Cloud Serving Benchmark (YCSB) defines a family of key--value workloads used
here as the common load profile: YCSB-A (50/50 read/update, Zipfian), YCSB-B (95/5
read/update), YCSB-C (read-only), YCSB-D (read-latest), YCSB-E (short range scans), and
YCSB-F (read--modify--write). Each algorithm profile declares an expected workload
(Chapter~\ref{ch:sota}), which the measurement driver uses for profile-aware routing
(Chapter~\ref{ch:methodology}).

\section{C++23 Metaprogramming}\label{sec:cpp23}

The engine relies on compile-time composition to avoid run-time dispatch in the hot path.
\emph{Concepts} constrain the building-block interfaces; the \emph{curiously recurring
template pattern} (CRTP) provides static polymorphism; \texttt{if~constexpr} and
\texttt{std::conditional\_t} select between a probe implementation and a cache-engine
default per axis; and C++23 \emph{modules} provide an ABI-stable boundary between the
permutation binaries and the builder. Together these facilities realize the
zero-cost abstraction on which the permutation matrix depends.
